    \documentclass[conference]{IEEEtran}
    \IEEEoverridecommandlockouts

    \usepackage[T1]{fontenc}
    \usepackage[utf8]{inputenc}
    \usepackage[ngerman]{babel}
    \usepackage{amsmath,amssymb,amsfonts}
    \usepackage{textcomp}
    \usepackage{xcolor}
    \usepackage{url}
    \usepackage{hyperref}

    \begin{document}

    \title{Diffie-Hellman Schlüsselaustausch\\
    {\footnotesize Gruppenarbeit IT-Sicherheit}}

    \author{\IEEEauthorblockN{}
    }

    \maketitle



    \section{Einleitung}

    Der Diffie-Hellman-Schlüsselaustausch (DH) wurde 1976 von Whitfield
    Diffie und Martin Hellman in ihrem Paper „New directions in
    cryptography``~\cite{dh1976} veröffentlicht. Vor dieser Arbeit war
    es ein offenes Problem, wie zwei Parteien über einen unsicheren Kanal
    einen gemeinsamen Schlüssel aushandeln können, ohne sich vorher
    persönlich getroffen zu haben.

    Das Verfahren löst dieses Problem und ist bis heute Bestandteil von
    vielen wichtigen Protokollen, zum Beispiel TLS, SSH oder IPsec. Beide 
    Parteien tauschen Werte über
    einen Kanal aus, den ein Angreifer mitlesen kann, und kommen am Ende
    trotzdem zu einem gemeinsamen Geheimnis, das der Angreifer nicht
    kennt. Möglich wird das durch eine mathematische Einwegfunktion, also
    eine Operation, die in eine Richtung einfach, in die andere Richtung
    aber praktisch nicht umkehrbar ist.

    In diesem Abschnitt wird zuerst die mathematische Grundlage erklärt,
    danach der Protokollablauf Schritt für Schritt durchgegangen und
    zum Schluss werden die Sicherheitsannahmen besprochen, auf denen DH
    beruht.

    \section{Mathematische Grundlagen}

    \subsection{Die Gruppe $\mathbb{Z}_p^*$}

    DH operiert in einer endlichen, zyklischen Gruppe. Klassischerweise
    ist das die multiplikative Gruppe der ganzen Zahlen modulo einer
    Primzahl~$p$:

    \begin{equation}
        \mathbb{Z}_p^* = \{1, 2, \ldots, p-1\}
    \end{equation}

    Diese Gruppe enthält alle Zahlen von 1 bis $p-1$. Die Gruppenoperation
    ist die Multiplikation modulo $p$, das heißt, das Ergebnis jeder
    Multiplikation wird durch $p$ geteilt und nur der Rest behalten.
    Dass $p$ eine Primzahl ist, ist wichtig, weil dadurch jedes Element
    der Gruppe ein multiplikatives Inverses hat. Wäre $p$ keine Primzahl,
    gäbe es Nullteiler und die Struktur würde nicht mehr funktionieren.

    \subsection{Generatoren}

    Ein \emph{Generator} (oder primitives Element) ist ein Element $g \in
    \mathbb{Z}_p^*$, das durch wiederholtes Potenzieren alle Elemente der
    Gruppe erzeugt:

    \begin{equation}
        \langle g \rangle = \{g^0, g^1, g^2, \ldots, g^{p-2}\} \bmod p
        = \mathbb{Z}_p^*
    \end{equation}

    Das heißt, wenn man $g$ immer wieder mit sich selbst multipliziert
    (und jedes Mal modulo $p$ rechnet), bekommt man irgendwann jede Zahl
    zwischen 1 und $p-1$ als Ergebnis. Nach $p-1$ Schritten landet man
    wieder bei 1, und der Zyklus beginnt von vorne. Die Anzahl der
    Elemente in $\langle g \rangle$ wird als \emph{Ordnung} von~$g$
    bezeichnet.

    Nicht jedes Element ist ein Generator. Zur Verdeutlichung ein
    Beispiel in $\mathbb{Z}_7^* = \{1, 2, 3, 4, 5, 6\}$. Für $g = 3$
    ergibt sich:
    \begin{align*}
    3^0 &= 1 \\
    3^1 &= 3 \\
    3^2 &= 9 \equiv 2 \pmod 7 \\
    3^3 &= 27 \equiv 6 \pmod 7 \\
    3^4 &= 81 \equiv 4 \pmod 7 \\
    3^5 &= 243 \equiv 5 \pmod 7 \\
    3^6 &= 729 \equiv 1 \pmod 7
    \end{align*}

    Damit erzeugt $g = 3$ alle Elemente von $\mathbb{Z}_7^*$ und ist
    folglich ein Generator. Für $g = 2$ gilt dagegen $2^3 \equiv 1 \pmod
    7$, das heißt, $2$ erzeugt nur die Untergruppe $\{1, 2, 4\}$ und ist
    kein Generator. Welche Elemente Generatoren sind, hängt also von $p$
    ab. In der Praxis wird $g$ zusammen mit $p$ in Standards festgelegt,
    zum Beispiel in den MODP-Gruppen aus RFC~3526~\cite{rfc3526}.

    \subsection{Das Diskrete-Logarithmus-Problem}

    Die Sicherheit von DH basiert auf folgendem Problem. Gegeben sind
    $g$, $p$ und ein Wert $y = g^x \bmod p$. Gesucht ist der Exponent
    $x$. Dieses Problem nennt man das \emph{Diskrete-Logarithmus-Problem}
    (DLP). Formal sucht man

    \begin{equation}
        x = \log_g y \pmod p.
    \end{equation}

    Für kleine $p$ ist das einfach, man probiert einfach alle Werte
    durch. Für große $p$ (in der Praxis mindestens 2048 Bit) gibt es
    aber keinen bekannten klassischen Algorithmus, der das in vernünftiger
    Zeit löst. Das Potenzieren in eine Richtung ($x \to g^x \bmod p$) ist
    also einfach, der Rückweg aber nicht. Genau diese Asymmetrie zwischen
    einfacher und schwerer Richtung ist das, was DH überhaupt erst
    möglich macht~\cite{cfrg-dh}. Alice hat dabei den Wissensvorteil,
    weil sie ihren Exponenten kennt, während ein Angreifer ihn nur aus
    dem öffentlichen Wert zurückrechnen müsste.

    \section{Protokollablauf}

    Das Protokoll wird hier am Beispiel von zwei Parteien, Alice und Bob,
    beschrieben. Es lässt sich problemlos auf drei oder mehr Parteien
    erweitern, was im Demo-Teil der Arbeit gezeigt wird.

    \subsection{Schritt 1: Öffentliche Parameter}

    Alice und Bob einigen sich zuerst auf eine große Primzahl~$p$ und
    einen Generator~$g$ mit $1 < g < p$. Diese Werte sind \emph{nicht}
    geheim. Sie können öffentlich bekannt sein und werden in der Praxis
    oft aus Standards übernommen, statt selbst generiert zu werden.

    \subsection{Schritt 2: Private Schlüssel}

    Alice wählt eine zufällige Zahl
    \begin{equation}
        a \in \{1, \ldots, p-1\}
    \end{equation}
    und behält diese geheim. Bob macht dasselbe und wählt sich ein
    zufälliges
    \begin{equation}
        b \in \{1, \ldots, p-1\}.
    \end{equation}
    Die Werte $a$ und $b$ sind die jeweiligen privaten Schlüssel. Sie
    verlassen niemals die Geräte ihrer Besitzer.

    \subsection{Schritt 3: Öffentliche Schlüssel}

    Alice berechnet ihren öffentlichen Schlüssel
    \begin{equation}
        A = g^a \bmod p
    \end{equation}
    und schickt diesen an Bob. Bob berechnet entsprechend
    \begin{equation}
        B = g^b \bmod p
    \end{equation}
    und schickt $B$ an Alice. Auch wenn ein Angreifer $A$ und $B$
    mitliest, kann er daraus wegen des DLP nicht auf $a$ oder $b$
    zurückschließen.

    \subsection{Schritt 4: Gemeinsames Geheimnis}

    Jetzt kann jede Partei das gemeinsame Geheimnis berechnen. Alice
    rechnet
    \begin{equation}
        k = B^a \bmod p
    \end{equation}
    und Bob rechnet
    \begin{equation}
        k' = A^b \bmod p.
    \end{equation}

    Dass beide auf denselben Wert kommen, sieht man durch Einsetzen:

    \begin{equation}
        k \equiv (g^b)^a \equiv g^{ab} \equiv (g^a)^b \equiv k' \pmod p
    \end{equation}

    Beide Berechnungen liefern $g^{ab} \bmod p$, weil bei der modularen
    Exponentiation $(g^a)^b = (g^b)^a = g^{ab}$ gilt. Das funktioniert,
    weil die Operation kommutativ ist. Der ausgehandelte Wert $k$ wird
    dann üblicherweise durch eine Schlüsselableitungsfunktion (Key
    Derivation Function, KDF) in einen symmetrischen Schlüssel für ein
    Verfahren wie AES überführt.

    \section{Sicherheitsannahmen}

    Die Sicherheit von DH beruht auf der Annahme, dass bestimmte Probleme
    in $\mathbb{Z}_p^*$ schwer zu lösen sind. Dabei werden zwei eng
    verwandte Probleme unterschieden.

    \subsection{Diskretes-Logarithmus-Problem (DLP)}

    Wie schon erklärt, lässt sich aus $A = g^a \bmod p$ der Wert $a$
    nicht effizient berechnen, solange $p$ groß genug ist. Würde es einen
    effizienten Algorithmus für das DLP geben, wäre DH komplett gebrochen,
    weil ein Angreifer aus den öffentlichen Schlüsseln direkt die
    privaten Schlüssel rekonstruieren könnte.

    \subsection{Computational Diffie-Hellman (CDH)}

    Das CDH-Problem ist eine etwas andere Formulierung. Gegeben sind
    $g$, $p$, $A = g^a \bmod p$ und $B = g^b \bmod p$. Gesucht ist
    $g^{ab} \bmod p$, also das gemeinsame Geheimnis. Das CDH-Problem ist
    höchstens so schwer wie das DLP. Wer das DLP lösen kann, kann auch
    CDH lösen. Ob die Umkehrung gilt, ist offen.

    Für DH reicht es aus, dass CDH schwer ist. Dass ein Angreifer also
    $A$ und $B$ kennt, hilft ihm nichts, solange er $g^{ab}$ nicht
    ausrechnen kann.

    \subsection{Wahl der Parameter}

    In der Praxis sind die Parameter entscheidend. Wählt man $p$ oder $g$
    schlecht, kann das Protokoll trotz korrekter Mathematik unsicher sein.
    Übliche Vorgaben~\cite{bsi-tr}:

    \begin{itemize}
        \item $p$ sollte mindestens 2048 Bit groß sein, besser 3072 Bit
            oder mehr.

        \item $p$ sollte eine \emph{sichere Primzahl} sein, also von der
            Form $p = 2q + 1$, wobei $q$ ebenfalls prim ist (sogenannte
            Sophie-Germain-Primzahl). Dadurch hat die relevante
            Untergruppe Primzahlordnung~$q$, was bestimmte
            Angriffsklassen wie den Pohlig-Hellman-Algorithmus
            verhindert.

        \item $g$ sollte ein Generator einer Untergruppe von Primzahlordnung
            sein.

        \item Die privaten Schlüssel $a$ und $b$ müssen kryptografisch
            sicher zufällig gewählt werden. Ein vorhersehbarer
            Zufallsgenerator würde alles aushebeln.
    \end{itemize}

    Die in RFC~3526~\cite{rfc3526} definierten MODP-Gruppen werden in der
    Praxis häufig verwendet, weil sie diese Eigenschaften erfüllen und
    ausreichend geprüft sind.

    \section{Elliptische-Kurven-Variante (ECDH)}

    Eine in der Praxis weit verbreitete Variante ist der Elliptic Curve
    Diffie-Hellman-Austausch (ECDH). Das Prinzip ist dasselbe, nur dass
    nicht in $\mathbb{Z}_p^*$ gerechnet wird, sondern in der Gruppe der
    Punkte einer elliptischen Kurve über einem endlichen Körper.

    Der Vorteil ist, dass das diskrete Logarithmusproblem auf elliptischen
    Kurven schwerer ist als in $\mathbb{Z}_p^*$. Dadurch reichen deutlich
    kürzere Schlüssel für ein vergleichbares Sicherheitsniveau. Während
    klassisches DH heute mindestens 2048-Bit-Schlüssel braucht, reichen
    bei ECDH 256 Bit aus~\cite{rfc8446}. Das ist gerade für mobile Geräte
    oder eingebettete Systeme relevant, weil weniger Rechenleistung und
    weniger Bandbreite benötigt werden.

    In modernen Protokollen wie TLS 1.3 ist ECDH inzwischen der
    Standardfall. Klassisches DH wird aus Performance- und
    Sicherheitsgründen kaum noch genutzt.

    \section{Ausblick: Post-Quanten-Kryptographie}

    Sowohl klassisches DH als auch ECDH beruhen auf der Schwierigkeit des
    diskreten Logarithmusproblems. Mit einem ausreichend großen
    Quantencomputer ließe sich dieses Problem allerdings durch den
    Shor-Algorithmus effizient lösen, was DH (und auch RSA) brechen
    würde~\cite{nist-pqc}.

    Aus diesem Grund wird aktuell an Post-Quanten-Kryptographie geforscht,
    also an Verfahren, die auch gegen Quantenangriffe sicher sein sollen.
    Dazu gehören unter anderem gitterbasierte Verfahren wie ML-KEM (früher
    Kyber), code-basierte Verfahren wie McEliece und hash-basierte
    Signaturen wie XMSS oder LMS. In der Praxis werden zunehmend
    hybride Verfahren empfohlen, die klassisches DH oder ECDH mit einem
    Post-Quanten-Verfahren kombinieren.

    \section{Zusammenfassung}

    DH löst das Problem des Schlüsselaustauschs über einen unsicheren
    Kanal. Beide Parteien wählen einen privaten Schlüssel, berechnen
    daraus einen öffentlichen Schlüssel und tauschen diesen aus. Durch
    die Eigenschaft $(g^a)^b = (g^b)^a$ erhalten beide am Ende denselben
    Wert $g^{ab} \bmod p$, ohne dass dieser je übertragen wurde.

    Die Sicherheit beruht darauf, dass das DLP und das verwandte
    CDH-Problem für genügend große, sauber gewählte Parameter praktisch
    nicht lösbar sind. In der hier vorgestellten Form schützt das
    Protokoll aber nur gegen passive Angreifer. Ein aktiver Angreifer,
    der sich in die Kommunikation einklinkt, kann DH ohne zusätzliche
    Maßnahmen brechen. Das ist Thema des nächsten Abschnitts (Demo und
    Man-in-the-Middle-Angriff).

    \section{Vermeidung von Angriffen}

    Wie bei allen kryptografischen Protokollen ist es auch bei DH wichtig, die richtigen Parameter zu wählen und zusätzliche Maßnahmen zu ergreifen, um Angriffe zu verhindern. Im Folgenden werden einige bekannte Angriffe und deren Gegenmaßnahmen beschrieben.

    \subsection{Man in the middle}
    Wie bereits erwähnt, kann ein aktiver Angreifer, der sich in die Kommunikation einklinkt,den DH-Schlüsselaustausch brechen. Um einen klassischen Man-in-the-Middle-Angriff zu verhindern, müssen die Parteien ihre öffentlichen Schlüssel authentifizieren. Ist erstmal sichergestellt, dass zwei Parteien tatsächlich miteinander kommunizieren, ist ein Angriff als Zwischenposten nicht mehr möglich. Das kann zum Beispiel durch digitale Signaturen oder durch die Verwendung von Zertifikaten geschehen. In TLS 1.3 wird dies durch die Verwendung von Zertifikaten und der Signatur der Handshake-Nachrichten erreicht.\cite[§4.4.2, §4.4.3]{rfc8446} Näheres dazu wird in Abschnitt xx erläutert. 

    \subsection{Logjam}
    Ein weiterer typischer Angriff ist der sogenannte Logjam-Angriff, bei dem ein Angreifer schwache Parameter (z.B. eine zu kleine Primzahl) erzwingt, um das DLP lösbar zu machen. Dabei manipuliert der Angreifer die Handshake-Nachrichten so, dass beide Parteien davon ausgehen, dass das System nur schwache Gruppen unterstützt, und dadurch auf unsichere Parameter zurückfallen. So werden dann zum Beispiel nur 512 Bit große Primzahlen verwendet statt die sichereren 2048 Bit.\cite{logjam} Moderne Protokolle wie TLS 1.3 haben solche schwachen Gruppen explizit ausgeschlossen, um dieses Problem zu vermeiden.\cite[§1.2]{rfc8446} Da aber viele Systeme noch TLS 1.2 und alte DHE\_EXPORT (absichtlich schwache Krypthografie) unterstützen, da diese nie aktiv entfernt wurden, ist es wichtig, auch hier auf die Wahl sicherer Parameter zu achten und schwache Gruppen aktiv zu deaktivieren.\cite{logjam, rfc7919}

    \subsection{Number Field Sieve}
    Da der DH-Schlüsselaustausch auf der Komplexität des DLP basiert, ist das Problem - und damit die Sicherheit des Systems - leicht lösbar sobald der Angreifer die genutzte Primzahl kennt. Es wurde empirisch gezeigt, dass eine einzige Vorberechnung für eine 1024-Bit-Primzahl ausreicht, um passiv 18\% der populären HTTPS-Sites, 66\% der IPsec-VPNs und 26\% der SSH-Server zu kompromittieren.\cite{logjam} Denn das ist der Punkt an dem der Number Field Sieve Algorithmus anssetzt. Dieser kann in einer sehr aufwändigen Vorberechnungsphase die Primzahl faktorisieren. Dabei sucht er sich viele kleine Primzahlen — die sogenannte Faktorbasis — und berechnet für jede davon den diskreten Logarithmus in der Gruppe modulo p. Dabei entsteht eine Art Nachschlagewerk pro Gruppe. In Phase 2 berechnet der Algorithmus dann konkrete Kombinationen der Faktorbasis um einen konkreten Wert A (bekannter Schlüssel) zu bekommen. Gelingt dies, so ist die zugehörige private Zahl a bekannt und das DLP gelöst.\cite{logjam} So kann im Anschluss jeder Server, der diese Primzahl verwendet, gebrochen werden. Welche Server diese Primzahl nutzt kann durch valide TLS-Handshakes oder durch Tools wie ZMap — ein 2013 entwickelter schneller Netzwerkscanner der das gesamte IPv4-Internet in unter 45 Minuten scannen kann\cite{zmap} - herausgefunden werden.  Daher ist es wichtig, dass Serverbetreiber sicherstellen, dass sie eine einzigartige Parametergruppe verwenden, um die Sicherheit ihrer Server zu gewährleisten.

    \section{Forward Secrecy}
    Ein wichtiger Aspekt moderner kryptografischer Protokolle ist die sogenannte Forward Secrecy (auch Perfect Forward Secrecy, PFS). Sie stellt sicher, dass selbst wenn ein Angreifer in der Zukunft den privaten Schlüssel eines Servers kompromittiert, vergangene Kommunikationssitzungen nicht entschlüsselt werden können. \cite[§1.2]{rfc8446}\cite{rfc4949} Das bedeutet, dass die Sicherheit der Kommunikation nicht von der langfristigen Sicherheit eines einzelnen Schlüssels abhängt. Ohne Forward Secrecy könnte ein Angreifer, der den privaten Schlüssel eines Servers stiehlt, alle früheren Aufzeichnungen der Kommunikation entschlüsseln, was ein erhebliches Risiko darstellt. 
    \\
    

    Beim DH-Schlüsselaustausch kann Forward Secrecy erreicht werden, indem für jede Sitzung ein neues Schlüsselpaar generiert wird. Diese Schlüssel werden auch Ephemeral Schlüssel genannt, weil sie nur für die Dauer einer Sitzung gültig sind. Sobald die Sitzung beendet ist, werden diese Schlüssel verworfen. Das bedeutet, dass selbst wenn der private Schlüssel eines Servers kompromittiert wird, die Schlüssel für vergangene Sitzungen nicht bekannt sind und daher nicht entschlüsselt werden können. In TLS 1.3 ist die Unterstützung von Forward Secrecy verpflichtend, was die Sicherheit der Kommunikation erheblich verbessert.\cite{logjam}\cite[§1.2]{rfc8446}
    Das selbe Prinzip kann auch bei der zuvor erwähnten Elliptischen-Kurven-Variante (ECDH) angewendet werden, um auch hier Forward Secrecy zu gewährleisten.
    \\
    
    
    Ein weiterer Punkt bei dem Forward Secrecy wichtig ist, ist die Post-Quanten-Kryptographie. Wie bereits erwähnt, können Quantencomputer bald in der Lage sein, das DLP zu lösen. Es wird von Sicherheitsbehörden wie der NSA und NIST angenommen, dass verschlüsselte Kommunikationen bereits heute aufgezeichnet werden, um sie später zu entschlüsseln.\cite{nist-pqc} Durch die Verwendung von Forward Secrecy wird dieses Risiko reduziert, da die Schlüssel für vergangene Sitzungen nicht mehr verfügbar sind, selbst wenn der langfristige private Schlüssel kompromittiert wird. So müsste ein Angreifer für jede Sitzung ein neues DLP lösen, was den Angriff erheblich erschwert.








    \begin{thebibliography}{00}

    \bibitem{dh1976}
    W. Diffie und M. Hellman, ``New directions in cryptography,'' 1976.
    [Online]. Verfügbar:
    \url{https://ee.stanford.edu/~hellman/publications/24.pdf}

    \bibitem{rfc3526}
    T. Kivinen und M. Kojo, ``More Modular Exponential (MODP)
    Diffie-Hellman groups for Internet Key Exchange (IKE),'' RFC 3526,
    IETF, Mai 2003. [Online]. Verfügbar:
    \url{https://www.rfc-editor.org/rfc/rfc3526}

    \bibitem{rfc8446}
    E. Rescorla, ``The Transport Layer Security (TLS) Protocol Version
    1.3,'' RFC 8446, IETF, Aug. 2018. [Online]. Verfügbar:
    \url{https://www.rfc-editor.org/rfc/rfc8446}

    \bibitem{cfrg-dh}
    E. Rescorla, ``Diffie-Hellman Key Agreement Method,'' RFC 2631,
    IETF, Juni 1999. [Online]. Verfügbar:
    \url{https://www.rfc-editor.org/rfc/rfc2631}

    \bibitem{bsi-tr}
    Bundesamt für Sicherheit in der Informationstechnik, ``Technische
    Richtlinie TR-02102-1: Kryptographische Verfahren, Empfehlungen und
    Schlüssellängen.'' [Online]. Verfügbar:
    \url{https://www.bsi.bund.de/SharedDocs/Downloads/DE/BSI/Publikationen/TechnischeRichtlinien/TR02102/BSI-TR-02102.html}

    \bibitem{nist-pqc}
    National Institute of Standards and Technology, ``Post-Quantum
    Cryptography Standardization.'' [Online]. Verfügbar:
    \url{https://csrc.nist.gov/projects/post-quantum-cryptography}

    \bibitem{logjam}
    D. Adrian et al., ``Imperfect Forward Secrecy: How Diffie-Hellman
    Fails in Practice,'' in \textit{Proc. 22nd ACM Conference on Computer
    and Communications Security (CCS'15)}, Denver, CO, Okt. 2015.
    [Online]. Verfügbar: \url{https://weakdh.org/imperfect-forward-secrecy-ccs15.pdf}

    \bibitem{rfc7919}
    D. Gillmor, ``Negotiated Finite Field Diffie-Hellman Ephemeral
    Parameters for TLS,'' RFC 7919, IETF, Aug. 2016.
    [Online]. Verfügbar: \url{https://www.rfc-editor.org/rfc/rfc7919}

    \bibitem{zmap}
    Z. Durumeric, E. Wustrow und J. A. Halderman, ``ZMap: Fast
    Internet-wide Scanning and Its Security Applications,'' in
    \textit{Proc. 22nd USENIX Security Symposium}, Washington, D.C.,
    Aug. 2013, S. 605--620.
    [Online]. Verfügbar:
    \url{https://www.usenix.org/system/files/conference/usenixsecurity13/sec13-paper_durumeric.pdf}

    \bibitem{rfc4949}
    R. Shirey, ``Internet Security Glossary, Version 2,''
    RFC 4949, IETF, Aug. 2007.
    [Online]. Verfügbar:
    \url{https://www.rfc-editor.org/rfc/rfc4949}

    \end{thebibliography}

    \end{document}